\section{Computation values, finite control and its boundaries}

\begin{annexitem}{annex:p17.meta.kernel}{thm:p17-meta-kernel}{The actual induced law and history sufficiency}
\begin{proof}
For fixed \(z,c\), each modeled instance has exactly one successor
summary under its specified executable transition. The sets
\(\{\omega:I(F_c(\omega,z))=z'\}\) are therefore disjoint and
partition the support of \(\mu_z\). Summing their nonnegative
weights gives nonnegative entries of \(K\), with total one.
At a successor of positive total weight, restrict the old law to
its partition cell and divide by that total. The result is a
probability law and is exactly the finite conditional law: for
each subset of that cell, its conditional mass times the cell's
mass recovers its old mass. If instances change, push this law
forward through the specified instance transition. At a cell of
mass zero no division is performed; a supplied continuation model
there has no effect on expectation from the current state.

Conditional on any observed history, the sufficiency hypothesis
identifies its next-summary distribution with the displayed kernel,
its feasible blocks with \(C(z)\), and its cost and stopping
payoff with the displayed state functions. Thus choosing a block
at that history gives the same one-step conditional expected
return as choosing it in state \(z\). Iterating conditional
finite sums gives the same statement for any continuation of
finite length; this is the input used by the backward proof below.
The hypothesis is not inferred merely from giving a summary a name.

For example, two histories may have retained respectively bit zero
and bit one, while a proposed summary erases which was retained.
A block returning that workspace bit then has deterministic output
zero at the first history and one at the second. No single next-state
kernel at the erased summary can equal both distributions. Retaining
the bit, or a sufficient conditional law for it, repairs this
particular violation. A compressed model that does not do so has
not satisfied the hypothesis.
\end{proof}
\end{annexitem}

\begin{annexitem}{annex:p17.meta.baseline}{thm:p17-meta-baseline}{Same-estimator improvement and the precise information-value bridge}
\begin{proof}
Because \(\alpha_c\) maximizes the new vector,
\(X_{\alpha_c}\geq X_\alpha\) at each outcome. Their difference
is therefore nonnegative. If \(\alpha_c=\alpha\), it is zero,
regardless of how the value of that action changed relative to
an earlier evaluation. Subtracting the declared common cost gives
the score \(-k_c\). The earlier-baseline score minus the
new-baseline score is
\[
 (E[\max_iX_i]-\widehat u_\alpha-k_c)
 -(E[\max_iX_i-X_\alpha]-k_c)
       =E[X_\alpha]-\widehat u_\alpha.
\]
This proves both the discrepancy formula and its exact zero condition.

For the calibrated case, let the computed output have finite
positive-probability outcomes \(y\), with masses \(q_y\),
and let \(u_i(\omega)\) be the intrinsic actual utility of
action \(i\). Set
\[
 \widehat u_i=\sum_\omega\mu(\omega)u_i(\omega),\qquad
 X_i(y)=\sum_\omega\Pr(\omega\mid y)u_i(\omega).
\]
Multiplying by \(q_y\) and summing over outputs recovers
\(\mu(\omega)\) for each instance, so
\(\sum_yq_yX_i(y)=\widehat u_i\). Zero-probability outputs
contribute nothing. This is the finite conditional-expectation
identity, proved here directly by regrouping a finite sum.
Optimal output-dependent external choice has value
\(\sum_yq_y\max_iX_i(y)\); the old optimal choice has value
\(\widehat u_\alpha=\max_i\widehat u_i\).
Their difference is exactly
\(E[\max_iX_i-X_\alpha]\).
It is the signalwise value formula of
\cref{thm:p17-information-policy-value}. That theorem's bound
applies because the old choice remains available and the output
can be ignored. A computation can implement such an information
structure, but its program, running time, workspace transition and
prediction law still have to be specified as in this chapter.

The coherence premise cannot be omitted. Suppose two estimates
are initially \((1,0)\) and afterwards \((3,2)\), deterministically.
The intention stays the first action. The new-baseline gross benefit
is zero; the earlier-baseline expression gives two. If actual
utilities did not change and the same external action is taken,
that two is a revision of evaluation, not an external improvement.
Using an imperfect evaluator does not justify calling it calibrated.
\end{proof}
\end{annexitem}

\begin{annexitem}{annex:p17.meta.subtree}{thm:p17-meta-subtree}{Both one-component formulas and the dependence counterexample}
\begin{proof}
First suppose only the inferior action \(\beta\) changes.
Every other estimate is fixed, and their maximum is the old
intention's value \(a\). After the computation the maximum is
\(\max(a,X)\), while the old intention still has value \(a\).
Their difference is \(\max(a,X)-a=(X-a)_+\).
Subtracting the common cost and taking expectation gives the
first formula, including the equality case \(X=a\) and any
tie rule.

If only the old intention changes, the other estimates have fixed
maximum \(b\). The new maximum is \(\max(X,b)\), and the
same later evaluator gives the old intention value \(X\).
Their difference is \(\max(X,b)-X=(b-X)_+\), giving the
second formula. A low revised value of the old intention can
therefore justify avoiding it, although the newly chosen value
is below an earlier overestimate. For a density, split the
expectation at the threshold. Below \(a\) the first positive
part is zero; above \(b\) the second is zero. The remaining
integrals are precisely those in the statement. Integrability
ensures they are finite.

If only one external action exists, it is also the action chosen
after every complete computation. The new-estimator difference
is identically zero by the preceding baseline proof; no maximum
over an empty set of alternatives is taken.

For the dependence example, take the first action as the fixed
old intention. In the joint law concentrated on \((0,0)\) and
\((2,2)\), the maximum always equals its first coordinate,
so the gross benefit is zero. Independent coordinates with the
same marginals have four equally likely outcomes. The maximum
is zero once and two three times, giving expectation \(3/2\).
The first coordinate's expectation is one, so the gross benefit
is \(1/2\). Thus marginal laws alone do not determine the
multiple-component formula. No probabilistic independence was
needed in either one-component derivation.
\end{proof}
\end{annexitem}

\begin{annexitem}{annex:p17.meta.bellman}{thm:p17-meta-bellman}{Existence, policy law and every backward-control case}
\begin{proof}
Every allowed policy has a finite tree of observed summaries and
choices, of depth at most \(N\), with termination whenever stop
is selected or the budget becomes zero. The probability of a
trace is the product of its policy choice probabilities and its
transition probabilities. Their outgoing masses sum to one at
each internal node. Induction on the finite depth therefore makes
the terminal trace masses sum to one, and defines every expected
return as a finite sum. A zero-probability history uses its
supplied continuation tables, without conditional division by
its probability from the initial state.

We prove optimality by induction on remaining budget. At zero,
the sole allowed choice is the implemented terminal action, so
every policy has value \(g(z)=V_0(z)\).
Suppose the conclusion holds with \(n-1\) blocks. At a history
ending in \(z\), a policy may first stop, giving \(g(z)\),
or choose \(c\), pay \(k(z,c)\), and observe \(z'\).
Conditional on that successor its remaining-policy value is at
most \(V_{n-1}(z')\) by induction and state sufficiency.
Thus that first computation has value at most \(Q_n(z,c)\).
Randomizing the first choice forms a convex combination of these
numbers and the stop payoff; this cannot exceed their finite maximum.

Choose an option attaining that maximum. If it is stop, stop.
If it is a computation, select it and use an induction-optimal
policy with budget \(n-1\) at each successor. This produces
value exactly \(Q_n(z,c)\). Selecting the first maximizing
option in a fixed finite ordering yields a deterministic
state/budget policy attaining \(V_n(z)\) simultaneously from
every state and every supplied history. An empty computation
set leaves only stop. Equality comparisons permit stop because
it already attains the maximum. A strict computation advantage
excludes stop from the maximizers.

Only finitely many states, budgets, options and successor entries
are involved. Rational sums, products and comparisons compute
every displayed value and choice in finitely many operations.
This establishes computability from supplied rational tables,
not a bound on the cost of obtaining them or of doing the
metalevel arithmetic online. Resource exclusions enter through
the allowed option lists before this induction; an excluded
block is never used in the constructed policy.
\end{proof}
\end{annexitem}

\begin{annexitem}{annex:p17.meta.policy-law}{thm:p17-meta-policy-law}{Grounded self-prediction and selected-action conditional means}
\begin{proof}
At zero remaining budget the controller stops, with expected
payoff \(g(z)\). At positive budget, its first choice stops
with probability \(p_n(z)\) and continues otherwise. Conditional
on continuation it chooses \(c\) with probabilities \(d_n\),
pays its cost and reaches the stated successor law. Conditioning
on these first choices gives exactly the displayed recursion for
\(W_n\). Applying the same argument to each successor expands
it into a sum over all complete traces. The products are the
trace probabilities, terminal contributions are the terminal
external payoffs, and every incurred cost is counted once at
the node where it occurs. This proves the claimed expectation
identity by induction on the finite depth. Optimality of \(V_n\)
against this particular policy gives \(W_n\leq V_n\).

The events \(\{A=a\}\) partition the finite outcomes. On
each one the realized utility \(U_A\) is \(U_a\).
Therefore
\[
 E[U_A]=\sum_a E[\mathbf1_{\{A=a\}}U_a]
       =\sum_{a:\Pr(A=a)>0}\Pr(A=a)E[U_a\mid A=a].
\]
Zero-probability events contribute zero and require no conditional
mean. In the example the larger-utility action always has realized
payoff one, so the left side is one. Each action is selected half
the time and each has unconditional mean \(1/2\). Replacing
the two conditional means, which are one, by those unconditional
means gives \(1/2\), proving failure of that substitution.
Conditioning must include the terminal clock as well when utility
depends on delay; the selected trace determines that clock.
\end{proof}
\end{annexitem}

\begin{annexitem}{annex:p17.meta.approximate}{thm:p17-meta-approximate}{The actual greedy-error induction and threshold certificate}
\begin{proof}
Treat stop as one of the finite options. Its exact
optimal-continuation score is \(g(z)\); a computation's is
\(Q_n(z,c)\). Write \(q(a)\) for these scores at a fixed
state/budget, let \(a^*\) maximize \(q\), and let \(\tilde a\)
maximize the approximations. The uniform score bound gives
\[
 q(\tilde a)\geq\widetilde q(\tilde a)-\epsilon_n
 \geq\widetilde q(a^*)-\epsilon_n
 \geq q(a^*)-2\epsilon_n.
\]
If the chosen option stops, its realized continuation loss is zero;
the displayed inequality bounds its total loss by \(2\epsilon_n\).
If it computes, then relative to its exact optimal-continuation
score, its subsequent approximate policy loses at most the
largest loss at budget \(n-1\), averaged with a probability
distribution. With \(L_0=0\), a uniform loss bound therefore
satisfies \(L_n\leq2\epsilon_n+L_{n-1}\).
Induction gives \(L_n\leq2\sum_{j=1}^n\epsilon_j\).
Actual policy loss is nonnegative by exact optimality.

For two options each approximated within \(\epsilon_n\),
their true difference is at least their approximate difference
minus \(2\epsilon_n\). A strict approximate advantage larger
than that preserves the true ordering. For a one-step net value
\(b\) approximated by \(\tilde b\) within \(\epsilon\),
one instead has \(b\geq\tilde b-\epsilon\).
Thus \(\tilde b>\epsilon\) proves \(b>0\); if all true
one-step values are nonpositive, such an estimate contradicts
the assumed uniform bound. These are different certificates:
the latter checks one-step improvement, while the former policy
bound requires accurate optimal-continuation scores at each
remaining budget. An empirical prediction table does not acquire
either uniform bound merely by being used as a score.
\end{proof}
\end{annexitem}

\begin{annexitem}{annex:p17.meta.regress}{thm:p17-meta-regress}{Nonuniqueness, positive loop cost and priced control}
\begin{proof}
For a real scalar \(g\), \(w=\max(g,w)\) holds if and only
if \(w\geq g\), so it has all these solutions. If instead
\(w=\max(g,w-c)\), with \(c>0\), its right side is at
least \(g\), so \(w\geq g\). Were \(w>g\), the maximum
would have to be \(w-c\), giving \(w=w-c\), impossible.
Thus \(w=g\), which indeed solves the equation because
\(g-c<g\). These examples distinguish a grounded value from
an unrestricted self-referential valuation constraint.

Let a controller's modeled policy return, including its selected
computations' direct costs but excluding the additional controller
charge, be \(J\). Its action-independent charge \(r\) makes
the return \(J-r\). Immediate action without invoking that
controller returns \(g\). Subtracting gives improvement
\(J-g-r\), positive exactly when \(J-g>r\), zero at equality.
No assumption that the controller selects an optimal computation
is needed for this comparison. In the numerical case,
\(J-g=1-3/4=1/4\), and
\(1/4-3/10=-1/20\), as stated.

One could insert a further action that evaluates whether to invoke
the controller. That inserts another transition, cost and choice;
it cannot make its own cost vanish. A finite hierarchy, an
explicit budget, or a compiled policy that acts without a fresh
optimization at some level can terminate this particular regress.
The scalar ungrounded equation shows why calling every such level
optimal does not establish a uniquely defined or implementable
procedure.
\end{proof}
\end{annexitem}

\begin{annexitem}{annex:p17.meta.two-block}{thm:p17-meta-two-block}{Every program outcome, budget value and overhead comparison}
\begin{proof}
The four equally weighted inputs are \((0,0),(0,1),(1,0),(1,1)\).
The deterministic two-block program returns evaluation four on
the first and last, and minus two on the other two. Hence the
initial utility mean for \(B\) is
\((4-2-2+4)/4=1\), exceeding the zero payoff of \(A\).
With no computation the implemented intention \(B\) returns
one in expectation.

After \(c_1\) returns either zero or one, its conditional
instance set contains two equally weighted inputs, one with each
parity. The mean for \(B\) remains \((4-2)/2=1\), and the
intention remains \(B\). If action follows immediately, the
terminal expected payoff is one, so including its cost gives
\(1-2/5=3/5\). Thus a one-block initial policy strictly prefers
stopping. This proves the one-step decision using its actual
internal-state transition, rather than assigning the partial
block an unexplained information benefit.

At either first-block state, executing \(c_2\) gives the
accurate terminal evaluation. On the equal-bit instance take
\(B\) and receive four; on the unequal-bit instance take \(A\)
and receive zero. Terminal expected payoff is therefore two.
Paying the second-block cost gives \(2-2/5=8/5\), which is
strictly above the stopping payoff one. There are no other
computational choices. With two blocks initially, its only
computation option thus gives
\(8/5-2/5=6/5\), strictly above one. The exact finite recursion
selects both blocks and the stated terminal decision. The four
trace returns, including both block costs, are respectively
\(16/5,-4/5,-4/5,16/5\); their average is \(6/5\).

The single-step approximation scores \(c_1\) negatively because
it commits after that block, while the actual two-block policy
retains the bit and continues. Its improvement is \(6/5-1=1/5\).
A compiled controller with extra total charge \(1/10\) returns
\(6/5-1/10=11/10\), still exceeding one. Its instruction
sequence, conditional terminal branch and costs are explicitly
specified, so this is an executable finite policy in the declared
architecture. A deadline admitting only one block excludes its
second block; the one-block comparison already proved that stop
is then optimal. No policy can evade that limit by invoking a
computation outside the declared feasible instruction interface.
\end{proof}
\end{annexitem}
